\chapter{The Workbench: Architecture and Axes}

\section{Design philosophy}

Rather than committing to a single architecture, the workbench exposes
every meaningful design choice as a switchable configuration parameter.
This makes fair comparison possible by construction: when one axis
changes, everything else is held fixed, so any observed difference in
AUROC is attributable to that choice alone. The entire architecture is
specified by a Hydra configuration tree, where each axis corresponds to
a named key in a structured YAML hierarchy. Any run in the experiment
database can be regenerated from a single command-line invocation.

\section{The base transformer}

The base model follows the standard transformer encoder pipeline:
particle tokens are embedded into a fixed-dimensional space, passed
through a stack of attention-plus-FFN blocks, aggregated into a single
event-level vector by a pooling operation, and classified by a head
network. This skeleton is fixed across all experiments; the axes
control what happens inside each stage.

Each particle token, a raw kinematic vector plus an optional
particle-type identifier, is projected into the model dimension $d$
by a linear embedding layer. The encoder consists of $L$ identical
blocks, each containing a multi-head self-attention layer followed by
a position-wise feed-forward network, with residual connections and
normalisation at positions controlled by axes A1 and A2. The pooled
vector is passed to a classifier head that produces the final logit.

\section{Axis groups}

The axes are organised into ten groups. Groups G and D define the
experimental conditions and data treatment. Groups T, E, P, A, F, B,
and C define the architecture. Group H sets model capacity. Three
cross-cutting shared parameter blocks (§K, §M, §S) are consumed by
multiple modules simultaneously and are documented at the end of this
section. Groups §R and §L cover training protocol and logging
respectively and are not architectural axes.

% ──────────────────────────────────────────────
\subsection*{G \quad Study framing}

G axes define what is being classified and with which model family.
They are not architectural choices and are not varied in the ablation
programme; they set the experimental conditions under which all other
axes are evaluated.

\textbf{G01 -- Task type.} Selects the training loop: supervised
classification or anomaly detection (autoencoder, GAN-AE, diffusion-AE).
This thesis focuses exclusively on supervised classification;
the anomaly detection loops are available in the workbench but not
studied here.

\textbf{G02 -- Model family.} Selects the model architecture family:
transformer, MLP, BDT, or autoencoder. All Part II chapters use the
transformer family; the BDT appears only in Ch.\,8 as a surrogate
model for configuration prediction, not as a classifier.

\textbf{G03 -- Classification task.} Defines which processes are
treated as signal and background. Five processes are available:
four-top-quark production ($tttt$), associated production with a Higgs
boson ($ttH$), with a $W$ boson ($ttW$), with two $W$ bosons ($ttWW$),
and with a $Z$ boson ($ttZ$). Three task definitions are used in this
thesis: $tttt$ vs.\ combined background ($ttH + ttW + ttWW + ttZ$) as
the primary task, $tttt$ vs.\ $ttH$ as a harder secondary task, and
all-five multiclass as an exploratory task.

% ──────────────────────────────────────────────
\subsection*{D \quad Data treatment}

D axes control what the model sees before any architectural processing.

\textbf{D01 -- Feature set.} Selects which kinematic features are
included in each token. The default is the four-feature set
(E, Pt, $\eta$, $\phi$); energy can be omitted to test whether the
model relies on explicit $E$ or whether Pt, $\eta$, and $\phi$ already
provide sufficient kinematic information.

\textbf{D02 -- MET treatment.} Controls whether missing transverse
energy is excluded or included as an event-level global. Tests how much
the invisible neutrino system contributes beyond the particle tokens.

\textbf{D03 -- Token ordering.} Sets whether tokens are presented in
input order, sorted by decreasing Pt, or randomly shuffled per event.
Tests whether the model exploits ordering as an implicit positional
signal.

% ──────────────────────────────────────────────
\subsection*{T \quad Tokenizer}

T axes control how raw particle features are converted into token
embeddings.

\textbf{T1 -- Tokenizer family.} Selects between \texttt{raw}
(kinematics only), \texttt{identity} (kinematics plus an explicit
particle-type label), and \texttt{binned} (discrete token, no
continuous features). Tests whether explicit particle identity is
necessary or recoverable from kinematics alone.

\textbf{T1-a -- PID embedding mode.} When T1 = \texttt{identity},
selects how the particle-type label is encoded: \texttt{learned},
\texttt{one\_hot}, or \texttt{fixed\_random}. Tests whether particle
identity needs a learned geometry or only a categorical code.

\textbf{T1-b -- PID embedding dimension.} Sets the capacity of the
PID embedding to 8, 16, or 32 dimensions. Tests the trade-off between
representational capacity and overfitting for the type sub-embedding.

% ──────────────────────────────────────────────
\subsection*{E \quad Positional encoding}

E axes control how position information is injected into the token
stream.

\textbf{E1 -- PE type.} Selects among \texttt{none},
\texttt{sinusoidal}, \texttt{learned}, and \texttt{rotary}. Tests
whether injecting position information helps a set-based classifier at
all.

\textbf{E1-a -- PE space.} Controls whether the positional signal is
applied before or after the input projection. Prerequisite: E1
$\in$ \{sinusoidal, learned\}.

\textbf{E1-a1 -- PE dimension mask.} Allows the positional encoding to
act selectively on a subset of feature dimensions. Prerequisite:
E1-a = \texttt{token}.

% ──────────────────────────────────────────────
\subsection*{P \quad Pre-encoder modules}

P axes add optional modules applied to raw particle features before the
main transformer encoder. Both require T1 $\in$ \{raw, identity\}.

\textbf{P1 -- Nodewise mass patch.} Optionally enriches each token with
invariant-mass features computed from its $k$-nearest neighbours in
Lorentz space, following Li et al. Tests whether explicit local mass
structure before the encoder improves token representations. Requires
D01 to include energy (index 0).

\textbf{P2 -- MIA pre-encoder.} Optionally inserts MIParT-style
interaction-attention blocks before the main encoder, where attention
is driven by pairwise interaction features rather than query-key
similarity. Tests whether physics-driven message passing as a
pre-encoder stage adds information beyond standard self-attention.

\textbf{P2-a -- MIA placement.} Controls whether MIA blocks are
prepended or appended relative to the main encoder. Prerequisite:
P2 = \texttt{true}.

% ──────────────────────────────────────────────
\subsection*{A \quad Attention and encoder block}

A axes control the internal design of each encoder block.

\textbf{A1 -- Normalisation policy.} Selects pre-norm, post-norm, or
NormFormer placement of layer normalisation within each block. Tests
whether normalisation position affects training stability for
particle-level inputs.

\textbf{A2 -- Normalisation type.} Chooses between LayerNorm and
RMSNorm. Tests whether removing the mean-subtraction step improves
efficiency without degrading performance.

\textbf{A3 -- Attention type.} Selects standard multi-head attention
or differential attention, which computes the difference of two softmax
maps to suppress diffuse attention patterns. Tests whether this
mechanism transfers from NLP to particle classification.

\textbf{A3-a -- Differential bias mode.} When A3 =
\texttt{differential}, controls how additive physics biases interact
with the two attention branches: \texttt{none}, \texttt{shared}, or
\texttt{split}. Prerequisite: A3 = \texttt{differential}.

\textbf{A4 -- Attention-internal normalisation.} Optionally applies
LayerNorm or RMSNorm per head after the weighted sum, before head
merging. Tests whether per-head normalisation stabilises attention
output scales. Applies to both standard and differential attention.

\textbf{A5 -- Causal masking.} Optionally applies an autoregressive
mask to the attention matrix. Tests whether imposing sequential
structure on an inherently unordered set of particles hurts, helps,
or is neutral.

% ──────────────────────────────────────────────
\subsection*{F \quad Encoder FFN realization}

F axes control the feed-forward network inside each encoder block. The
effective realization follows a priority rule: MoE $>$ KAN $>$ standard.

\textbf{F1 -- FFN realization.} Selects between a standard two-layer
MLP, a Kolmogorov--Arnold Network (KAN), and a Mixture-of-Experts (MoE)
FFN. Tests whether learnable spline activations or conditional
computation improve per-token transformations.

\textbf{F1-a -- KAN FFN variant.} When F1 = \texttt{kan}, selects
among \texttt{hybrid}, \texttt{bottleneck}, and \texttt{pure}
architectures. Prerequisite: F1 = \texttt{kan}.

\textbf{F1-b -- MoE scope.} Controls whether expert routing applies
to encoder FFNs, the classifier head, or both. Prerequisite:
F1 = \texttt{moe}.

% ──────────────────────────────────────────────
\subsection*{B \quad Physics-informed attention biases}

B axes inject physics structure as additive biases on the attention
logits. All families require T1 $\in$ \{raw, identity\} and share a
zero-initialised scalar gate, so adding any bias leaves the baseline
unchanged at initialisation. Multiple families can be active
simultaneously using \texttt{+} notation.

\textbf{B1 -- Bias activation set.} The master switch selecting which
bias families are active: \texttt{lorentz\_scalar},
\texttt{typepair\_kinematic}, \texttt{sm\_interaction},
\texttt{global\_conditioned}, or combinations thereof.

\textbf{B1-L1 -- Lorentz feature set.} Selects the subset of pairwise
Lorentz-invariant features fed into the bias network, from the minimal
set \{m$^2$, $\Delta R$\} through the MIParT basis to the full feature
set. Active when \texttt{lorentz\_scalar} $\in$ B1.

\textbf{B1-L2 -- Lorentz MLP type.} Selects whether the network
mapping Lorentz features to bias scores uses a standard MLP or a KAN.
Active when \texttt{lorentz\_scalar} $\in$ B1.

\textbf{B1-L5 -- Lorentz sparse gating.} Adds per-feature sigmoid
gates providing a direct measure of learned feature importance. Active
when \texttt{lorentz\_scalar} $\in$ B1.

\textbf{B1-T2 -- Type-pair initialisation.} Controls how the learnable
$N_\text{type} \times N_\text{type}$ interaction table is initialised:
from scratch, from a binary SM mask, or from fixed SM coupling
constants. Active when \texttt{typepair\_kinematic} $\in$ B1.

\textbf{B1-T3 -- Type-pair freeze.} Controls whether the type-pair
table is frozen after initialisation or remains trainable. Tests
whether the table is better used as a fixed physics prior or a
learnable parameter. Active when \texttt{typepair\_kinematic} $\in$ B1.

\textbf{B1-S1 -- SM interaction mode.} Selects the fidelity of the SM
interaction prior: \texttt{binary}, \texttt{fixed\_coupling}, or
\texttt{running\_coupling}. Active when \texttt{sm\_interaction}
$\in$ B1.

\textbf{B1-G1 -- Global-conditioned mode.} Controls how MET modulates
pairwise attention biases: as a global scale factor or as a directional
signal. Requires D02 = \texttt{true}. Active when
\texttt{global\_conditioned} $\in$ B1.

% ──────────────────────────────────────────────
\subsection*{C \quad Classifier head}

C axes control the final mapping from the pooled event representation
to class logits.

\textbf{C1 -- Head realization.} Selects between a \texttt{linear}
head, a \texttt{kan} head, and a \texttt{moe} head. Tests whether
specialised nonlinearities or conditional computation at the decision
boundary improve discrimination.

\textbf{C2 -- Pooling strategy.} Selects how the token sequence is
aggregated into a single event-level vector: \texttt{cls} token,
\texttt{mean} pooling, or \texttt{max} pooling. Placed here because
pooling is logically part of the readout stage rather than the encoder
block.

% ──────────────────────────────────────────────
\subsection*{H \quad Model-size hyperparameters}

H axes set model capacity and are not architectural design choices in
the sense of the other groups. They are tracked separately and are the
primary axis of variation in Ch.\,9. The convenience key
\texttt{model\_size\_key} bundles dim ($d$), depth ($L$), heads, and
FFN hidden dim into named presets spanning roughly 50k to 4M
parameters.

% ──────────────────────────────────────────────
\subsection*{Cross-cutting blocks: §K, §M, §S}

Three shared parameter blocks are consumed simultaneously by multiple
modules. §K (KAN hyperparameters: grid size, spline order, grid range)
affects every active KAN module, including the encoder FFN, classifier
head, and Lorentz bias network. §M (MoE hyperparameters: number of
experts, top-k, routing level) affects all active MoE sites. §S (shared
pairwise backbone) computes the pairwise feature matrix $F_{ij}$ once
per forward pass and shares it across all active B1 families and the
MIA pre-encoder.

\section{Experimental methodology}

Each axis is studied in isolation: one axis is varied across its
settings while all others are held at their defaults, forming an
orthogonal sweep. This ensures observed differences are attributable to
the axis under study and keeps the number of required experiments
linear in the number of axes rather than exponential.

Every result reported in this thesis is averaged over at least three
independently seeded runs, with error bars representing the standard
deviation across seeds. The area under the ROC curve (AUROC) is the
primary evaluation metric throughout: it measures the probability that
the classifier assigns a higher score to a signal event than to a
background event, and is threshold-independent. Signal efficiency at
fixed background rejection is reported as a complementary metric in
selected chapters.

All runs are tracked in a Weights \& Biases project where every axis
setting, hyperparameter, and evaluation metric is logged under a
structured key scheme mirroring the axis reference directly.
